\begin{abstract}
La scarsità di dati etichettati e lo sbilanciamento delle classi rappresentano ostacoli critici nello sviluppo di modelli di Deep Learning per la diagnostica medica. Questo lavoro esplora l'efficacia delle Generative Adversarial Networks (GAN) come strumento di data augmentation per migliorare la classificazione della polmonite in immagini radiografiche pediatriche \cite{kermany2018identifying}. Utilizzando un dataset intrinsecamente sbilanciato, proponiamo l’impiego di una Conditional Deep Convolutional Wasserstein GAN con Gradient Penalty (cDC-WGAN-GP), sintetizzando i contributi fondamentali di \cite{mirza2014conditionalgenerativeadversarialnets}, \cite{radford2015unsupervised}, \cite{arjovsky2017wasserstein} e \cite{gulrajani2017improved}. Una ResNet18 \cite{He_2016_CVPR} viene utilizzata come classificatore baseline e successivamente addestrata sul dataset aumentato. I risultati indicano che l'integrazione di campioni sintetici della classe minoritaria permette di ottimizzare l'F1-score, riducendo la polarizzazione del modello. [...continue DA...].
\end{abstract}